\recipe{Ian's Gingered Beef}
\label{recipe:gingered_beef}
\makes{6 portions}
\worktime{\SI{35} \minute}
\totaltime{\SI{3} \hour}
\source{Ian Smith}

\begin{ingredients}

  \item Rapeseed oil 
  \item \qty{1.25}{\kilo \gram} beef shin or brisket cut into very large chunks 
  \item 2 onions 
  \item \SI{50}{\gram} ginger 
  \item 3 garlic cloves 
  \item small bunch coriander 
  \item \SI{2}{\teaspoon} Chinese five-spice powder 
  \item 6 whole star anise 
  \item \SI{1}{\teaspoon} black peppercorn 
  \item \SI{100}{\gram} dark brown muscovado sugar 
  \item \SI{50}{\milli \litre} light soy sauce 
  \item \SI{50}{\milli \litre} dark soy sauce 
  \item \SI{2}{\tablespoon} tomato purée 
  \item beef stock 
\ingroup{To serve}
  \item thumb-sized chunk ginger shredded into matchsticks 
  \item \SI{1}{\tablespoon} vegetable or sunflower oil 
  \item Cooked jasmine rice
\end{ingredients}


\begin{method}
  	Heat a little of the oil in a large flameproof dish. Add the beef
chunks, in batches, and fry until browned. When each batch is
browned, transfer the beef to another dish. Very roughly chop the
onions, ginger, garlic and coriander stalks. Put in a food processor
and whizz to a paste. 

Wipe any oil out of the dish you browned the beef in. Add the paste
with a good splash of water and gently fry, scraping up any beef bits,
until the paste is fragrant and softened (add more water if the paste
sticks). Stir in the five-spice, star anise and peppercorns, cook for 1
min, then add the sugar, soy sauces and tomato purée. Return the beef
and any juices to the dish, then stir in enough stock to just about
cover. Bring to a gentle simmer. Heat oven to 160C/140C fan/gas 3.
Cover the dish, put in the oven and cook for 2½ hrs until the beef is
really tender.

Lift the beef out of the sauce into a dish, to keep warm. Boil the
sauce until reduced by about half and thickened. Meanwhile, fry the
ginger in the oil until golden and crispy. Return the beef to the sauce.
serve the beef spooned over rice and scattered with the crispy ginger.
  
\end{method}



\begin{figure}[h]
  	\centering
  	
	\begin{tikzpicture}
  		% Save the image in a box to get its dimensions
  		\sbox0{\includegraphics[width = \linewidth]{images/MissingImage.png}}%
  
  		% Draw a clipping rectangle with rounded corners
  		\path[clip, rounded corners=0.5cm] (0,0) rectangle (\wd0,4cm);
  
  		% Place the image inside the clipped area
  		\node[inner sep=0pt, outer sep=0pt, anchor=center] at (0.5\linewidth, 4cm) {\usebox0};
	\end{tikzpicture}

\end{figure}


